\section*{Capítulo \ref{ch:b2:diatomic-mos} --- Orbitales moleculares de las moléculas diatómicas}

\begin{solution}{exo:b2:diatomic-mos:1}
$E_+ = (-13.6 - 9.0)/1.60 = \qty{-14.1}{eV}$; $E_- = (-13.6 + 9.0)/0.40 = \qty{-11.5}{eV}$.
Dos electrones en $\psi_+$: $-28.25$ frente a $2 \times (-13.6) = \qty{-27.2}{eV}$, ligados por
\qty{1.05}{eV} en este modelo rudimentario.
\end{solution}

\begin{solution}{exo:b2:diatomic-mos:2}
\ce{He2}: $(2 - 2)/2 = 0$; \ce{He2+}: $(2 - 1)/2 = 0.5$; \ce{Li2}: 1; \ce{B2}: seis electrones de
valencia, $\sigma_{2s}^2\sigma_{2s}^{*2}\pi_{2p}^2$, \omterm{def:b2:diatomic-mos:bond-order}{orden de enlace} 1.
\end{solution}

\begin{solution}{exo:b2:diatomic-mos:3}
\ce{B2} (dos electrones desapareados en los dos orbitales $\pi_{2p}$, con mezcla s--p) y
\ce{O2} (dos en $\pi^*$). \ce{C2}, \ce{N2} y \ce{F2} son \omterm{def:b2:diatomic-mos:magnetism}{diamagnéticas}.
\end{solution}

\begin{solution}{exo:b2:diatomic-mos:4}
No nulo: (1s, $2\mathrm p_z$), ($2\mathrm p_x$, $2\mathrm p_x$), (2s, $2\mathrm p_z$).
Nulo por simetría: (1s, $2\mathrm p_x$) y ($2\mathrm p_x$, $2\mathrm p_y$).
\end{solution}

\begin{solution}{exo:b2:diatomic-mos:5}
La ionización arranca un electrón enlazante: \omterm{def:b2:diatomic-mos:bond-order}{orden de enlace} 2.5 en lugar de 3, luego un enlace más largo y
más débil, como muestra $D_0(\ce{N2+}) = \qty{8.71}{eV} < D_0(\ce{N2}) = \qty{9.76}{eV}$.
\end{solution}

\begin{solution}{exo:b2:diatomic-mos:6}
Los orbitales del oxígeno están más bajos (el oxígeno es más electronegativo): por el
\cref{thm:b2:diatomic-mos:heteronuclear}, los \omterm{def:b2:diatomic-mos:bonding}{orbitales enlazantes} tienen más peso sobre el
átomo más bajo, el oxígeno, y los orbitales más altos, sobre el carbono. El orbital ocupado
más alto, un orbital $\sigma$ de carácter mayoritariamente carbonado que apunta en dirección opuesta al oxígeno, es
el par no enlazante del carbono.
\end{solution}

\begin{solution}{exo:b2:diatomic-mos:7}
Solo el $2\mathrm p_z$ del flúor (a lo largo del eje) tiene la simetría del 1s del hidrógeno:
forman un orbital $\sigma$ con más peso sobre F y un $\sigma^*$ con más peso sobre H. Los
$2\mathrm p_x$, $2\mathrm p_y$ del flúor (solapamiento nulo) y su 2s, muy bajo (demasiado lejos en energía),
permanecen no enlazantes: tres pares no enlazantes. Los ocho electrones de valencia llenan el 2s, el $\sigma$
y los dos p no enlazantes.
\end{solution}

\begin{solution}{exo:b2:diatomic-mos:8}
$\bar\alpha = \qty{-12}{eV}$, $\Delta = \qty{4}{eV}$, $\sqrt{4 + 9} = \qty{3.61}{eV}$:
$E = \qty{-15.6}{eV}$ y \qty{-8.4}{eV}. Para el nivel más bajo, $(\alpha_B - E)c_B = -\beta
c_A$: $1.61c_B = 3.0c_A$, $c_A/c_B = 0.535$, $c_B^2 = 1/(1 + 0.286) = 0.78$.
\end{solution}

\begin{solution}{exo:b2:diatomic-mos:9}
Ocho electrones de valencia: $\sigma_{2s}^2\sigma_{2s}^{*2}\pi_{2p}^4$, con el orbital $\sigma_{2p}$
(empujado por encima de $\pi_{2p}$ por la mezcla) vacío. \omterm{def:b2:diatomic-mos:bond-order}{Orden de enlace} $(6 - 2)/2 = 2$, debido a los
dos orbitales $\pi$ llenos.
\end{solution}

\begin{solution}{exo:b2:diatomic-mos:10}
NO, once electrones de valencia, uno en $\pi^*$: \omterm{def:b2:diatomic-mos:bond-order}{orden de enlace} 2.5. \ce{NO+}, diez: \omterm{def:b2:diatomic-mos:bond-order}{orden de enlace}
3. $D_0(\ce{NO+}) = 6.51 + 13.62 - 9.26 = \qty{10.87}{eV}$: mucho más fuerte que NO, ya que
el electrón arrancado era antienlazante.
\end{solution}

\begin{solution}{exo:b2:diatomic-mos:11}
$E_+ + E_- = \dfrac{(\alpha + \beta)(1 - S) + (\alpha - \beta)(1 + S)}{1 - S^2} = \dfrac{2(\alpha -
\beta S)}{1 - S^2}$. Supera $2\alpha$ cuando $\alpha - \beta S > \alpha - \alpha S^2$, es decir,
$\beta < \alpha S$ (dividiendo por $-S < 0$). Con cuatro electrones, dos en cada orbital, la
energía total está por encima de la de los orbitales separados: dos orbitales llenos se repelen.
\end{solution}

\begin{solution}{exo:b2:diatomic-mos:12}
\omterm{def:b2:diatomic-mos:bond-order}{Órdenes de enlace} 2.5, 2, 1.5 y 1: en este orden los enlaces se alargan y se debilitan. El peróxido
de hidrógeno contiene el enlace sencillo O--O del peróxido (unidad \ce{O2^2-}); \ce{KO2} contiene
el ion superóxido \ce{O2-}, \omterm{def:b2:diatomic-mos:magnetism}{paramagnético}.
\end{solution}

\begin{solution}{pb:b2:diatomic-mos:1}
\textbf{1.} Doce.
\textbf{2.} $\sigma_{2s}^2\,\sigma_{2s}^{*2}\,\sigma_{2p}^2\,\pi_{2p}^4\,\pi_{2p}^{*2}$, con los dos electrones
$\pi^*$ en orbitales distintos.
\textbf{3.} $(8 - 4)/2 = 2$.
\textbf{4.} Dos: \ce{O2} es \omterm{def:b2:diatomic-mos:magnetism}{paramagnético}.
\textbf{5.} Aparea todos los electrones, cuando dos de ellos están solos en dos orbitales
$\pi^*$ degenerados.
\textbf{6.} $D_0 = 2 \times 246.79 = \qty{493.6}{kJ/mol}$, \qty{5.12}{eV}.
\textbf{7.} \ce{O2+}: $\pi^{*1}$, \omterm{def:b2:diatomic-mos:bond-order}{orden de enlace} 2.5.
\textbf{8.} \ce{O2-}: $\pi^{*3}$, \omterm{def:b2:diatomic-mos:bond-order}{orden de enlace} 1.5.
\textbf{9.} \ce{O2^2-}: $\pi^{*4}$, \omterm{def:b2:diatomic-mos:bond-order}{orden de enlace} 1.
\textbf{10.} \ce{O2+}, uno; \ce{O2-}, uno; \ce{O2^2-}, ninguno.
\textbf{11.} \ce{O2+} $<$ \ce{O2} $<$ \ce{O2-} $<$ \ce{O2^2-}.
\textbf{12.} El ion peróxido.
\textbf{13.} \ce{O2+}, \ce{O2} y \ce{O2-}.
\textbf{14.} De un orbital $\pi^*$.
\textbf{15.} El orbital $\pi^*$, antienlazante, está por encima del nivel 2p del átomo: su
electrón está menos ligado.
\textbf{16.} Dos caminos de \ce{O2} a $\ce{O} + \ce{O+} + \mathrm e^-$: disociar y luego
ionizar un átomo, $D_0(\ce{O2}) + I(\ce{O})$; o ionizar la molécula y luego disociar el ion,
$I(\ce{O2}) + D_0(\ce{O2+})$. Son iguales: $D_0(\ce{O2+}) = 5.12 + 13.62 - 12.07 = \qty{6.66}{eV}$.
\textbf{17.} Mayor que $D_0(\ce{O2}) = \qty{5.12}{eV}$: arrancar un electrón antienlazante
refuerza el enlace (orden 2.5).
\textbf{18.} $D_0(\ce{N2}) = 2 \times 470.82/96.485 = \qty{9.76}{eV}$; $D_0(\ce{N2+}) = 9.76 +
14.53 - 15.58 = \qty{8.71}{eV}$.
\textbf{19.} El orbital ocupado más alto de \ce{N2} es enlazante, por debajo del nivel 2p del
átomo.
\textbf{20.} Cuando el electrón arrancado procede de un \omterm{def:b2:diatomic-mos:bonding}{orbital antienlazante}.
\textbf{21.} La regla de Hund (número máximo de espines paralelos en orbitales degenerados).
\textbf{22.} Ambos electrones $\pi^*$ en el mismo orbital, con los espines apareados.
\textbf{23.} Dos electrones en el mismo orbital se repelen más, y pierden la
estabilización de canje de los espines paralelos.
\textbf{24.} El \omterm{def:b2:diatomic-mos:bond-order}{orden de enlace} sigue siendo 2; todos los electrones apareados: \omterm{def:b2:diatomic-mos:magnetism}{diamagnético}.
\textbf{25.} \ce{O2+} tiene un \omterm{def:b2:diatomic-mos:bond-order}{orden de enlace} $\boldsymbol{2.5}$ y una energía de disociación de
$\boldsymbol{\approx \qty{6.66}{eV}}$, frente a 2 y \qty{5.12}{eV} para \ce{O2}.
\end{solution}
